\section{Introduction}
\label{sec:introduction}

A self-adjoint operator, a prescribed mean counting law, and a nonzero-time
periodic-orbit fluctuation are three different mathematical achievements.
The first gives a real discrete spectrum.  The second controls its averaged
density.  Only the third begins to expose geometry beyond the mean.  Keeping
these levels separate is especially important in Hilbert--P\'olya-motivated
model building, where a well-designed Weyl law can otherwise be mistaken for
an arithmetic spectral mechanism.

This paper studies the local fluctuation problem for an explicit pair of
semiclassical Schr\"odinger operators
\begin{equation}
 P_{a,\h}=-\frac{\h^2}{2}\Delta+
 2\pi\exp\!\left(\pi|\Psi_a(q)|^2\right),
 \qquad
 P_{0,\h}=-\frac{\h^2}{2}\Delta+2\pi e^{\pi|q|^2},
 \label{eq:intro-pair}
\end{equation}
where \(\Psi_a\) is a centered, determinant-one H\'enon automorphism and the
flagship parameter is \(a=1.02\).  The configuration change preserves area,
so the two classical Hamiltonians have the same exact sublevel volume.  That
identity fixes a common mean clock but does not force the quantum spectra to
coincide.  The published deterministic-chaos study supplies the broader
prime-distribution motivation \citep{Wang2026PrimeChaos}, while the specific
area-preserving H\'enon/operator construction and the value \(a=1.02\) were
introduced in the subsequent H\'enon preprint
\citep{Wang2026HenonPreprint}.  The present work asks a
narrower question: can one isolate, with proofs and auditable numerical
certificates, a nonzero-time relative spectral term belonging to a specific
classical orbit?

The difficulty is not the formal Gutzwiller expression.  Rigorous
fixed-energy trace formulas are established
\citep{DuistermaatGuillemin1975,BrummelhuisUribe1991,PaulUribe1995,CombescureRalstonRobert1999}.
The model-specific work is to verify that every stationary point allowed by
the chosen time cutoff has been accounted for.  A competing warped orbit, a
radial family, a repeated return, or a transverse degeneracy would change
the coefficient or invalidate the isolated-orbit formula.  An inserted
microlocal observable would make isolation easier, but it would also make the
left-hand side depend on eigenfunctions.  Our target is instead a finite-rank
relative trace determined only by the two eigenvalue lists.

The central conclusion has two deliberately separate parts.  Analytically,
for each fixed sufficiently small positive energy excess
\(\delta=E-2\pi\), a blow-up of the complete energy shell leaves one
primitive fast Lyapunov return in the positive-time window and no radial
return.  This produces an eigenvalue-only relative Gutzwiller coefficient.
Quantitatively, a computer-assisted proof certifies the same distinguished
local branch and its transverse determinant on
\(0\le\epsilon\le0.101\), where \(\delta=\epsilon^2\).  The second result
does not quantify the small-energy threshold in the first because the root
complement and global return cover remain incomplete.

Our contributions are as follows.
\begin{enumerate}[leftmargin=*,label=\textup{(\roman*)}]
 \item We express the energy-localized relative propagator trace as an
 ordinary finite sum and as an exact Stieltjes functional of the relative
 counting staircase.  No first-resolvent trace-class assumption is needed.
 \item We compute the bottom normal form of the H\'enon well, identify the
 fast Lyapunov family, and derive its limiting period, action slope,
 transverse determinant, and first nonlinear period coefficient.  At
 \(a=1.02\),
 \[
 T_+^0=0.6638439766792985,\qquad
 D_+^0=3.8627220445155035.
 \]
 \item We prove that there is a nonquantitative \(\delta_{\rm tr}>0\) such
 that, for every fixed \(0<\delta<\delta_{\rm tr}\), one primitive fast
 orbit supplies the positive-time relative trace term
 \[
 i\widehat g(T_+)\frac{T_+}{2\pi\sqrt{|\det(I-P_+)|}}
 e^{iS_+/\h}+O(\h).
 \]
 The absence of an observable makes this coefficient recoverable from the
 two spectra alone.
 \item We certify one connected real-analytic primitive full-return branch,
 unique inside the frozen local boxes for
 \(0\le\epsilon\le0.101\), by 51 primary slabs and 50 guarded bridges at
 each of 128 and 256 MPFR bits.  The proof archive contains 202 successful
 validated-flow/Krawczyk jobs and 202 exact-rational Krawczyk replays.
 \item On that local branch, we prove the exact reduction
 \(\det(I-D\Pi_\epsilon)=4-\operatorname{tr}M_\epsilon\) without assuming
 semisimplicity of the unit multipliers, and validate the uniform inequality
 \(\det(I-D\Pi_\epsilon)>3\).
\end{enumerate}

The numerical spectral calculation at \(\delta=0.01\) is reported for a
different purpose.  Its finest normalized value,
\(1.0065230645+0.0133004473i\), is consistent with the analytic amplitude
and phase convention.  It is not evidence that
\(0.01<\delta_{\rm tr}\), because the latter inequality has not been
proved.  We retain the larger pre-asymptotic deviations in the same frozen
ladder rather than fitting them away.

\begin{figure}[t]
\centering
\setlength{\fboxsep}{7pt}
\begin{minipage}{0.96\textwidth}
\centering
\fbox{\parbox{0.88\textwidth}{\centering
\textbf{Spectral input:} two energy-localized eigenvalue lists\\
\(\Downarrow\) exact finite-rank/Stieltjes identities}}
\vspace{2mm}

\fbox{\parbox{0.88\textwidth}{\centering
\textbf{Analytic spine:} bottom normal form + complete-shell blow-up\\
\(0<\delta<\delta_{\rm tr}\), with \(\delta_{\rm tr}>0\) nonquantitative\\
\(\Downarrow\) one positive-time relative Gutzwiller coefficient}}
\vspace{2mm}

\fbox{\parbox{0.88\textwidth}{\centering
\textbf{Quantitative sidecar:} A4.12 local-box branch + A4.13 determinant gap\\
\(0\le\epsilon\le0.101\), \(\delta=\epsilon^2\)\\
\(\dashrightarrow\) R401-SC at \(\delta=0.01\) is diagnostic only}}
\vspace{2mm}

\fbox{\parbox{0.88\textwidth}{\centering
\textbf{Open arrows:} local complement \(\to\) phase/global cover
\(\to\) quantitative \(\delta_{\rm tr}\); independently,
endogenous prime times \(\to\) arithmetic explicit formula}}
\end{minipage}
\caption{Proof architecture and claim boundary.  Solid implications belong
to the nonquantitative analytic trace theorem or to the separate local-box
computer-assisted theorem.  Dashed implications remain open.  In
particular, the explicit branch interval does not license the trace formula
at \(\delta=0.01\), and neither route supplies an arithmetic prime trace.}
\label{fig:architecture}
\end{figure}

The evidence layers in \cref{fig:architecture} govern the language of the
paper.  We call an exact algebraic or functional-calculus statement an
identity, a hand proof a theorem or proposition, the validated local result
a computer-assisted theorem, and the frozen finite-dimensional computation
a numerical diagnostic.  The open arithmetic gate is not softened by any of
these labels.

The paper is organized as follows.  \Cref{sec:related} positions the result
against fixed-energy trace formulas, Lyapunov families, relative spectral
containers, and validated periodic-orbit continuation.
\Cref{sec:model-trace,sec:normal-form} build the exact spectral object and
the bottom oracle.  \Cref{sec:analytic-trace} proves the nonquantitative
one-orbit theorem.  \Cref{sec:certified-branch} states and analyzes the
explicit local-branch and monodromy certificates.
\Cref{sec:spectral-diagnostic} reports the frozen spectral diagnostic, and
\cref{sec:scope} records what is still absent from a Hilbert--P\'olya route.
The appendices give the proof and reproduction details.
